% ch19_monoidal_as_categorified_monoids.tex — Volume II Chapter 7

\chapter{Monoidal Categories as Categorified Monoids}

\begin{overviewbox}
Chapter~12 introduced monoidal categories and their coherence theorem. This
chapter takes a step back and asks: what \emph{is} a monoidal category from a
higher categorical perspective? The answer is that a monoidal category is a
\term{categorified monoid}: a monoid in which the underlying set has been
replaced by a category, and the equations of a monoid have been replaced by
natural isomorphisms.

This perspective makes the coherence theorem inevitable (the natural isos must
be compatible), clarifies the distinction between strict and weak monoidal
categories, and connects monoidal categories to the theory of monads through
the concept of a ``monoid in a monoidal category.'' We also develop monoidal
functors in their three flavors (lax, oplax, strong) and conclude with the
beautiful observation that monoids in various monoidal categories recover all
the algebraic structures one cares about.
\end{overviewbox}


% ═══════════════════════════════════════════════════════════════════════════
\begin{nameorigin}
\term{Categorification} and its inverse \term{decategorification} are
informal-but-precise terms popularized in the 1990s by \emph{Louis
Crane}, \emph{Igor Frenkel}, and \emph{John Baez}, in the context of
constructing categorical analogues of algebraic structures (especially
TQFTs and quantum groups). \emph{Categorify} = ``promote a set to a
category, equations to natural isos, etc.''; \emph{decategorify} =
``flatten back to underlying set / hom-set.'' The terminology has
become standard for the philosophical operation of moving between
$n$-categorical levels.

The \term{Eckmann-Hilton argument} (Eckmann-Hilton 1962) shows that two
compatible monoid structures on a set must coincide and be commutative.
Categorifying gives the $\mathbb{E}_2$/braided picture (Vol IV Ch.~56);
the 1-categorical original is foundational to monoidal-category theory.
\end{nameorigin}

\section{Decategorification}
\label{sec:decategorification}
% ═══════════════════════════════════════════════════════════════════════════

\subsection{Informally: Going Down One Level}

\term{Decategorification} means: take a categorical structure, forget about
morphisms (or collapse them), and see what algebraic structure remains. It is a
lossy map from the richer world to the poorer world, but it illuminates what
the richer world is a categorification \emph{of}.

For monoidal categories: if we take a monoidal category $(\mathcal{V}, \otimes,
I)$ and decategorify by collapsing all morphisms (keeping only the objects and
recording whether they are isomorphic), we obtain a \emph{commutative monoid}
when the monoidal category is symmetric --- and a \emph{monoid} in general.

\subsection{Expanding: The Decategorification Map}

Given a monoidal category $\mathcal{V}$, form the set $|\mathcal{V}|$ of
isomorphism classes of objects. The tensor product $\otimes$ descends: if
$A \cong A'$ and $B \cong B'$, then $A \otimes B \cong A' \otimes B'$. So
$\otimes$ induces a binary operation on $|\mathcal{V}|$, with unit $[I]$.

The associator $\alpha_{A,B,C} : (A \otimes B) \otimes C \cong A \otimes (B \otimes C)$
becomes an equality $([A][B])[C] = [A]([B][C])$ in $|\mathcal{V}|$. The unitors
become $[I][A] = [A]$ and $[A][I] = [A]$.

Result: $(|\mathcal{V}|, \cdot, [I])$ is a monoid.

\begin{dialog}
So a monoidal category is to a monoid as a category is to a set?

Exactly. A set is a category with only identity morphisms. A monoid is a set with
a binary operation. A monoidal category is a category with a tensor product
(binary operation on objects, functorial on morphisms). The equations of a monoid
become isomorphisms in the monoidal category. This is categorification.
\end{dialog}

\subsection{Formalizing: The Categorification Slogan}

\begin{center}
\begin{tabular}{ll}
\textbf{Monoid} & \textbf{Monoidal category} \\
\hline
Set of elements & Category of objects and morphisms \\
Binary operation $m \cdot n$ & Tensor functor $A \otimes B$ \\
Unit element $e$ & Unit object $I$ \\
Associativity equation & Associator natural isomorphism $\alpha$ \\
Unit equations & Unitor natural isomorphisms $\lambda, \rho$ \\
--- & Coherence: pentagon and triangle commute \\
\end{tabular}
\end{center}

The coherence diagrams (pentagon for $\alpha$, triangle for $\lambda, \rho$)
have no monoid analogue because in a monoid, the corresponding equations are
strict. In a monoidal category, the isomorphisms need to be ``glued together''
consistently --- that is what coherence demands.


% ═══════════════════════════════════════════════════════════════════════════
\section{Strict vs.\ Weak Monoidal Categories}
\label{sec:strict-vs-weak}
% ═══════════════════════════════════════════════════════════════════════════

\subsection{Informally: When Coherence Is Trivial}

A \term{strict monoidal category} is one where the associator and unitors are
\emph{identities}, not merely natural isomorphisms:
\begin{equation*}
  (A \otimes B) \otimes C = A \otimes (B \otimes C), \qquad
  I \otimes A = A = A \otimes I.
\end{equation*}
These are on-the-nose equalities of objects. In a strict monoidal category,
parenthesization is irrelevant by definition, not merely up to isomorphism.

The prototypical example is $(\mathbf{Set}, \times, \{*\})$ \emph{if} we choose
a specific product, but most constructions in nature produce only weak monoidal
categories with non-identity coherence isomorphisms.

\subsection{Expanding: Why Strict Is Rare}

Consider $(\mathbf{Vect}_k, \otimes_k, k)$: the tensor product of vector spaces.
We have $(U \otimes V) \otimes W \cong U \otimes (V \otimes W)$, but this is a
canonical isomorphism, not an equality. The two sides have the same dimension
but are \emph{different sets} of tensors. Similarly, $k \otimes V \cong V$ via
the scalar multiplication isomorphism, but $k \otimes V$ and $V$ are not
literally equal.

Strict monoidal categories do arise: the endofunctor category $[\mathcal{C},
\mathcal{C}]$ with composition is strict (functor composition is associative
and has strict units). The category of strict 2-functors between 2-categories
is also strict.

\subsection{Formalizing: The Definitions}

\begin{definition}
A \term{monoidal category} $(\mathcal{V}, \otimes, I, \alpha, \lambda, \rho)$
is called \term{strict} if $\alpha$, $\lambda$, $\rho$ are all identity natural
transformations (equivalently, identity morphisms at every object).
\end{definition}

\begin{definition}
A \term{weak monoidal category} (or simply monoidal category) has $\alpha$,
$\lambda$, $\rho$ as natural isomorphisms satisfying the pentagon and triangle
axioms.
\end{definition}

The pentagon axiom says the two ways to reassociate $(A \otimes B) \otimes (C
\otimes D)$ to $A \otimes (B \otimes (C \otimes D))$ agree. The triangle axiom
says $\rho_A \otimes \mathrm{id}_B = (\mathrm{id}_A \otimes \lambda_B) \circ
\alpha_{A,I,B}$.


% ═══════════════════════════════════════════════════════════════════════════
\section{Mac Lane's Coherence Theorem}
\label{sec:maclane-coherence}
% ═══════════════════════════════════════════════════════════════════════════

\subsection{Informally: Why Every Diagram Commutes}

When we write $A \otimes B \otimes C \otimes D$ in a monoidal category, there
are many ways to parenthesize this expression, and between any two
parenthesizations there are multiple composite natural isomorphisms built from
$\alpha$, $\lambda$, $\rho$. Mac Lane's coherence theorem guarantees they all
agree.

\begin{theorem}[Mac Lane's Coherence Theorem]
\label{thm:maclane-coherence}
In any monoidal category, every diagram built from instances of $\alpha$,
$\lambda$, $\rho$ and their inverses, using tensor product and composition,
commutes.
\end{theorem}

This is why we write $A \otimes B \otimes C$ without parentheses: any
parenthesization gives an isomorphic result, and the isomorphisms between
different parenthesizations are all equal.

\subsection{Expanding: The Normal Form Strategy}

The proof strategy is to reduce every parenthesization to a \term{normal form}
(fully right-associated, with no $I$ factors) using the coherence isomorphisms,
and then show that any two paths to the normal form give the same result.

The key lemma: any two composites of coherence isomorphisms between the same
source and target are equal. This is proved by induction on the complexity of
the composites, using the pentagon and triangle axioms as the base cases.

\subsection{Formalizing: Key Step via Prooflog}

We prove the key step: the pentagon axiom implies that any two ways to
reassociate a four-fold tensor product agree.

\begin{prooflog}
\proofstep{Let $P$ be the pentagon: $(\mathrm{id}_A \otimes \alpha_{B,C,D})
\circ \alpha_{A,B \otimes C, D} \circ (\alpha_{A,B,C} \otimes \mathrm{id}_D)
= \alpha_{A,B,C \otimes D} \circ \alpha_{A \otimes B, C, D}$.}{Pentagon axiom
for $(A,B,C,D)$.}
\proofstep{There are five distinct parenthesizations of $A \otimes B \otimes C
\otimes D$, corresponding to the five vertices of the Mac Lane pentagon.}{Vertices:
$((AB)C)D$, $(A(BC))D$, $A((BC)D)$, $A(B(CD))$, $(AB)(CD)$.}
\proofstep{The five edges of the pentagon are instances of $\alpha$. The pentagon
axiom says the diagram of these five edges commutes.}{Each edge applies $\alpha$
at one level.}
\proofstep{For a general diagram of coherence morphisms, we define the
\emph{normal form functor} $N$ that reduces any parenthesization to a unique
right-associated form $A_1 \otimes (A_2 \otimes (\cdots \otimes A_n)\cdots)$
using $\alpha$.}{Definition of $N$ by structural recursion.}
\proofstep{Any coherence morphism $\theta : P \to Q$ factors as $N(Q)^{-1} \circ
n_\theta \circ N(P)$ where $n_\theta$ is a morphism in normal form.}{By
induction on the structure of $\theta$.}
\proofstep{The pentagon axiom implies that $n_\theta$ is uniquely determined by
its source and target normal forms.}{The pentagon says there is only one way to
pass between adjacent parenthesizations in normal form.}
\proofstep{Therefore any two coherence morphisms with the same source and target
have the same normal form component, hence are equal.}{Two paths through the
same source/target reduce to the same normal-form morphism.}
\end{prooflog}

\subsection{Reorganizing: What Coherence Buys}

The coherence theorem has a slick reformulation:

\begin{corollary}
Every monoidal category is monoidally equivalent to a strict monoidal category.
\end{corollary}

\begin{proof}[Sketch]
Build the ``strictification'' $\mathcal{V}^{\mathrm{st}}$ whose objects are
\emph{finite lists} of objects of $\mathcal{V}$ and whose tensor product is
list concatenation (strict!). The functor $\mathcal{V}^{\mathrm{st}} \to
\mathcal{V}$ sends a list to its tensor product. The coherence theorem implies
this is an equivalence.
\end{proof}

This means we may always \emph{assume without loss of generality} that a
monoidal category is strict when doing abstract arguments. In concrete situations
(like $\mathbf{Vect}$) the non-strict form is necessary, but for proving
abstract theorems, strictification reduces to the easier case.


% ═══════════════════════════════════════════════════════════════════════════
\section{Monoidal Functors}
\label{sec:monoidal-functors}
% ═══════════════════════════════════════════════════════════════════════════

\subsection{Informally: Preserving the Tensor Structure}

A functor between monoidal categories should ``preserve the tensor product,''
but there are three flavors of how well it does so.

A \term{lax monoidal functor} $F : \mathcal{V} \to \mathcal{W}$ comes with a
\emph{comparison morphism}
\begin{equation*}
  \phi_{A,B} : F(A) \otimes F(B) \to F(A \otimes B)
\end{equation*}
and a unit comparison $\phi_0 : I_\mathcal{W} \to F(I_\mathcal{V})$, satisfying
coherence conditions. The morphism $\phi$ goes in the direction ``product of
images maps to image of product,'' but is not required to be an isomorphism.

An \term{oplax monoidal functor} has $\phi$ going the other direction:
$\phi_{A,B} : F(A \otimes B) \to F(A) \otimes F(B)$.

A \term{strong monoidal functor} (sometimes just ``monoidal functor'') has $\phi$
an isomorphism in both cases.

\subsection{Expanding: Examples}

\begin{example}{Free Monoid Is Lax}
The free monoid functor $L : \mathbf{Set} \to \mathbf{Mon}$ is lax monoidal
(with $\mathbf{Set}$ and $\mathbf{Mon}$ both having cartesian product as their
tensor). The comparison $L(A) \times L(B) \to L(A \times B)$ sends a pair of
words $(w, v)$ to the sequence of pairs. This is not an isomorphism (a word in
$L(A \times B)$ contains more interleaved information than a pair of words).
\end{example}

\begin{example}{Forgetful Functor Is Lax}
The forgetful functor from monoids to sets is lax monoidal in the reverse
direction: it is \emph{oplax}, with $|M \times N| \to |M| \times |N|$ the
identity (here it happens to be strict).
\end{example}

\begin{example}{Symmetric Monoidal Functors in Algebra}
A strong monoidal functor $F : \mathbf{Ab} \to \mathbf{Ab}$ (with tensor
product) that is additionally symmetric carries the structure of a Hopf algebra
or similar in some categorical sense.
\end{example}

\subsection{Formalizing: The Definitions}

\begin{definition}
A \term{lax monoidal functor} $(F, \phi, \phi_0) : (\mathcal{V}, \otimes, I)
\to (\mathcal{W}, \otimes, J)$ consists of:
\begin{itemize}
  \item A functor $F : \mathcal{V} \to \mathcal{W}$.
  \item A natural transformation $\phi_{A,B} : FA \otimes FB \to F(A \otimes B)$.
  \item A morphism $\phi_0 : J \to FI$.
\end{itemize}
These must satisfy the coherence conditions: two diagrams expressing compatibility
of $\phi$ with $\alpha$ and with $\lambda$, $\rho$.
\end{definition}

\begin{definition}
A \term{strong monoidal functor} is a lax monoidal functor where $\phi$ and
$\phi_0$ are isomorphisms.
\end{definition}

\begin{howtosay}
\begin{itemize}
  \item ``Lax monoidal'' means the comparison goes \emph{toward} the image:
        $FA \otimes FB \to F(A \otimes B)$.
  \item ``Oplax monoidal'' means the comparison goes \emph{away} from the
        image: $F(A \otimes B) \to FA \otimes FB$.
  \item ``Strong monoidal'' means both are isomorphisms (so it doesn't matter
        which direction you specify).
  \item In the literature, a ``monoidal functor'' with no adjective often means
        \emph{strong} monoidal. Check the convention.
\end{itemize}
\end{howtosay}

\subsection{Reorganizing: Functors and Algebras}

A key reason to study lax monoidal functors: a lax monoidal functor sends monoids
to monoids. If $M$ is a monoid in $\mathcal{V}$ (see Section~\ref{sec:monoids-in-monoidal})
and $F$ is lax monoidal, then $FM$ is a monoid in $\mathcal{W}$. The multiplication
of $FM$ is $FM \otimes FM \xrightarrow{\phi} F(M \otimes M) \xrightarrow{F\mu} FM$.

This is why lax monoidal functors are more useful than mere functors between
monoidal categories: they respect the algebraic structure.


% ═══════════════════════════════════════════════════════════════════════════
\section{Monoids in Monoidal Categories}
\label{sec:monoids-in-monoidal}
% ═══════════════════════════════════════════════════════════════════════════

\subsection{Informally: Algebra Inside a Category}

Once we have a monoidal category $(\mathcal{V}, \otimes, I)$, we can ask: what
does a ``monoid'' look like inside $\mathcal{V}$? The answer is an object $M$
with a multiplication $\mu : M \otimes M \to M$ and a unit $\eta : I \to M$,
satisfying associativity and unit laws in the form of commutative diagrams.

This is the correct abstraction because it recovers all the classical notions
of ``algebra'' for different choices of $\mathcal{V}$:

\begin{center}
\begin{tabular}{lll}
$\mathcal{V}$ & Monoidal structure & Monoid in $\mathcal{V}$ \\
\hline
$\mathbf{Set}$ & Cartesian $\times$ & Ordinary monoid \\
$\mathbf{Ab}$ & Tensor $\otimes_\mathbb{Z}$ & Ring \\
$\mathbf{Vect}_k$ & Tensor $\otimes_k$ & $k$-algebra \\
$\mathbf{Grp}$ & Direct product & Group (with extra structure) \\
$[\mathcal{C}, \mathcal{C}]$ & Composition & Monad \\
\end{tabular}
\end{center}

\subsection{Expanding: The Monad Connection}

The last row deserves emphasis: a \emph{monad} on $\mathcal{C}$ is precisely a
monoid in the monoidal category $[\mathcal{C}, \mathcal{C}]$ of endofunctors
with composition.

\begin{itemize}
  \item The object $M$ is an endofunctor $T : \mathcal{C} \to \mathcal{C}$.
  \item The multiplication $\mu : T \circ T \to T$ is the monad multiplication.
  \item The unit $\eta : \mathrm{Id} \to T$ is the monad unit.
  \item Associativity: $\mu \circ (T\mu) = \mu \circ (\mu T)$.
  \item Unit laws: $\mu \circ (T\eta) = \mathrm{id}_T = \mu \circ (\eta T)$.
\end{itemize}

These are exactly the monad axioms from Chapter~8. The monoidal structure of
$[\mathcal{C}, \mathcal{C}]$ is the reason monads are ``monoids in endofunctors.''

\subsection{Formalizing: Monoids Precisely}

\begin{definition}
Let $(\mathcal{V}, \otimes, I)$ be a monoidal category. A \term{monoid} in
$\mathcal{V}$ is a triple $(M, \mu, \eta)$ where $M \in \mathcal{V}$,
$\mu : M \otimes M \to M$, and $\eta : I \to M$, such that:
\begin{align*}
  \mu \circ (\mu \otimes \mathrm{id}_M) &= \mu \circ (\mathrm{id}_M \otimes \mu)
  \circ \alpha_{M,M,M},
  \quad \text{(associativity)} \\
  \mu \circ (\eta \otimes \mathrm{id}_M) &= \lambda_M,
  \quad \text{(left unit)} \\
  \mu \circ (\mathrm{id}_M \otimes \eta) &= \rho_M.
  \quad \text{(right unit)}
\end{align*}
A \term{monoid homomorphism} $(M, \mu, \eta) \to (M', \mu', \eta')$ is a
morphism $f : M \to M'$ with $f \circ \mu = \mu' \circ (f \otimes f)$ and
$f \circ \eta = \eta'$.
\end{definition}

\begin{definition}
A \term{comonoid} in $\mathcal{V}$ is a monoid in $\mathcal{V}^{\mathrm{op}}$:
an object $C$ with comultiplication $\delta : C \to C \otimes C$ and counit
$\varepsilon : C \to I$ satisfying co-associativity and co-unit laws.
\end{definition}

\subsection{Reorganizing: The Eckmann--Hilton Argument}

In a braided monoidal category (one with a coherent swap $\sigma_{A,B} : A
\otimes B \cong B \otimes A$), monoids with compatible braidings are
\emph{commutative}. The Eckmann--Hilton argument shows that a monoid in a
category of monoids must be commutative:

\begin{proposition}[Eckmann--Hilton]
If $M$ is simultaneously a monoid $(M, \mu, \eta)$ and $(M, \nu, \epsilon)$
in $\mathbf{Set}$ and the two operations are compatible ($\mu$ is a monoid
homomorphism for $\nu$), then $\mu = \nu$ and both are commutative.
\end{proposition}

Categorically, this says that a monoid in $(\mathbf{Mon}, \times, \{e\})$ is
the same as a commutative monoid. This is why double loop spaces are abelian
group objects: they carry two compatible multiplication structures that force
commutativity.

\subsection{Simplifying: The Internal Algebra Slogan}

\begin{equation*}
  \boxed{
    \text{Monoid in } (\mathcal{V}, \otimes, I)
    = \text{algebra internal to } \mathcal{V}
    \text{ with respect to } \otimes
  }
\end{equation*}

This single observation subsumes rings, algebras, and monads as special cases.
The monoidal structure of $\mathcal{V}$ determines what kind of algebra is
possible; different monoidal structures give different theories.


% ═══════════════════════════════════════════════════════════════════════════
\section*{Exercises}
\addcontentsline{toc}{section}{Exercises}
% ═══════════════════════════════════════════════════════════════════════════

\begin{exercisesbox}
\begin{qa}
\qaitem{What is decategorification of a monoidal category?}{Forget morphisms, keep isomorphism classes; the tensor product descends to a binary operation making $|\mathcal{V}|$ a monoid.}
\qaitem{What monoid does $(\mathbf{Vect}_k, \otimes_k, k)$ decategorify to?}{The monoid of isomorphism classes of vector spaces under $\otimes$, isomorphic to $(\mathbb{N}, +, 0)$ via dimension.}
\qaitem{What is a strict monoidal category?}{One where the associator and unitors are identities (not merely natural isomorphisms).}
\qaitem{Give an example of a strict monoidal category.}{The endofunctor category $[\mathcal{C},\mathcal{C}]$ with composition.}
\qaitem{Why is $(\mathbf{Vect}_k, \otimes_k, k)$ not strict?}{Because $(U \otimes V) \otimes W$ and $U \otimes (V \otimes W)$ are isomorphic but not literally equal as sets.}
\qaitem{State Mac Lane's coherence theorem.}{Every diagram built from $\alpha$, $\lambda$, $\rho$ and their inverses commutes.}
\qaitem{What is the key proof strategy for coherence?}{Reduce each parenthesization to a normal form (right-associated, $I$-free) and show any two composites agree by the pentagon and triangle axioms.}
\qaitem{What does the coherence theorem imply about strictification?}{Every monoidal category is monoidally equivalent to a strict one.}
\qaitem{What are the five vertices of the Mac Lane pentagon for $(A,B,C,D)$?}{$((AB)C)D$, $(A(BC))D$, $A((BC)D)$, $A(B(CD))$, $(AB)(CD)$.}
\qaitem{What is a lax monoidal functor?}{A functor with comparison morphisms $\phi_{A,B} : FA \otimes FB \to F(A\otimes B)$ and $\phi_0 : J \to FI$ satisfying coherence, but not required to be isomorphisms.}
\qaitem{What is an oplax monoidal functor?}{Same as lax but with comparisons going the other direction: $F(A\otimes B) \to FA \otimes FB$.}
\qaitem{What is a strong monoidal functor?}{A lax monoidal functor where all comparison morphisms are isomorphisms.}
\qaitem{Why do lax monoidal functors send monoids to monoids?}{Given $\mu: M\otimes M \to M$ and lax comparison $\phi$, compose $FM \otimes FM \xrightarrow{\phi} F(M\otimes M) \xrightarrow{F\mu} FM$.}
\qaitem{What is a monoid in a monoidal category $(\mathcal{V}, \otimes, I)$?}{An object $M$ with $\mu: M\otimes M \to M$ and $\eta: I \to M$ satisfying associativity and unit diagrams.}
\qaitem{What kind of structure is a monoid in $(\mathbf{Ab}, \otimes_\mathbb{Z}, \mathbb{Z})$?}{A ring (associative and unital).}
\qaitem{What kind of structure is a monoid in $([\mathcal{C},\mathcal{C}], \circ, \mathrm{Id})$?}{A monad on $\mathcal{C}$.}
\qaitem{State the Eckmann--Hilton argument.}{A set with two compatible monoid structures must have them equal and commutative.}
\qaitem{What does Eckmann--Hilton say about a monoid in $(\mathbf{Mon}, \times)$?}{It is a commutative monoid.}
\qaitem{What coherence data does a monoidal category have beyond the tensor functor?}{Associator $\alpha$, left unitor $\lambda$, right unitor $\rho$ satisfying pentagon and triangle axioms.}
\qaitem{What is a comonoid in $\mathcal{V}$?}{An object with comultiplication $C \to C\otimes C$ and counit $C \to I$ satisfying co-associativity and co-unit laws.}
\end{qa}
\end{exercisesbox}


% ═══════════════════════════════════════════════════════════════════════════
\section*{Chapter Summary}
\addcontentsline{toc}{section}{Chapter Summary}
% ═══════════════════════════════════════════════════════════════════════════

\begin{summarybox}
\textbf{Categorification.} A monoidal category is a categorified monoid:
equations become natural isomorphisms (coherence data). Decategorifying a
monoidal category recovers a monoid on isomorphism classes of objects.

\textbf{Strict vs.\ weak.} Strict monoidal categories have identity coherence
morphisms. Most natural examples are only weakly monoidal. By Mac Lane's
coherence theorem, every weak monoidal category is equivalent to a strict one,
so we can always strictify for abstract arguments.

\textbf{Monoidal functors.} Lax monoidal functors carry comparison morphisms
$FA \otimes FB \to F(A \otimes B)$ that need not be isomorphisms. They send
monoids to monoids. Oplax and strong are the directional variants.

\textbf{Monoids in monoidal categories.} An object $M$ with multiplication and
unit morphisms satisfying categorical associativity/unit laws. This subsumes:
ordinary monoids (in $\mathbf{Set}$), rings (in $\mathbf{Ab}$), $k$-algebras
(in $\mathbf{Vect}_k$), and monads (in $[\mathcal{C},\mathcal{C}]$). The Eckmann--Hilton
argument shows that double-monoid structures force commutativity.
\end{summarybox}


% ═══════════════════════════════════════════════════════════════════════════
\section*{Formal Restatement}
\addcontentsline{toc}{section}{Formal Restatement}
% ═══════════════════════════════════════════════════════════════════════════

\begin{formalrestatement}
\textbf{Monoidal category.} $(\mathcal{V}, \otimes, I, \alpha, \lambda, \rho)$
with $\otimes : \mathcal{V} \times \mathcal{V} \to \mathcal{V}$, natural isos
$\alpha_{A,B,C} : (A \otimes B) \otimes C \cong A \otimes (B \otimes C)$,
$\lambda_A : I \otimes A \cong A$, $\rho_A : A \otimes I \cong A$, satisfying
the pentagon and triangle axioms.

\medskip
\textbf{Coherence theorem.} Every diagram of coherence morphisms commutes.
Every monoidal category is monoidally equivalent to a strict one.

\medskip
\textbf{Lax monoidal functor.} $(F, \phi, \phi_0)$ with $\phi_{A,B} : FA
\otimes FB \to F(A \otimes B)$ and $\phi_0 : J \to FI$ natural and coherent.
Strong if $\phi, \phi_0$ invertible; oplax if reversed.

\medskip
\textbf{Monoid in $\mathcal{V}$.} $(M, \mu : M \otimes M \to M, \eta : I \to
M)$ with $\mu(\mu \otimes \mathrm{id}) = \mu(\mathrm{id} \otimes \mu)\alpha$
and $\mu(\eta \otimes \mathrm{id}) = \lambda$, $\mu(\mathrm{id} \otimes \eta)
= \rho$. Instances: monoids, rings, algebras, monads.

\medskip
\noindent\textit{Chapter~20 enriches this picture by replacing the hom-sets
of a category with objects of a monoidal category $\mathcal{V}$, yielding the
full theory of $\mathcal{V}$-enriched categories.}
\end{formalrestatement}
